%% =====================================================================
%%  CHAPTER 8: RISK ANALYSIS AND MITIGATION           (budget 2 pages)
%%  chapters/ch8-risk.tex
%% =====================================================================
\pcChapter{risk}{3}{ניתוח סיכונים ותוכנית הפחתה}
\begin{pcguide}[מה כותבים כאן]
סיכונים: מה עלול להשתבש ואיך תתמודדו. כמתים כל סיכון: \mth{R = P \times I} (הסתברות כפול השפעה, בסולם \rng{1}{5}).
סיכון קריטי: \mth{R \ge 12}, ולו חייבת להיות תוכנית הפחתה ותוכנית מגירה. כללו סיכונים מכל סוג, לא רק טכניים: נתונים (נתונים שלא הגיעו), אבטחה וסיווג, ארגוניים ולוח זמנים.
\end{pcguide}
\section{כימות הסיכונים}\label{sec:riskquant}
הסיכונים הוערכו כהערכת תכנון בסולם הסתברות \mth{P} והשפעה \mth{I} של \rng{1}{5}, וחומרתם \mth{R=P\times I}: אדום \mth{R\ge 12}, כתום
\rng{5}{11} וירוק \mth{R\le 4}. המספרים אינם תדירויות שנמדדו בארגון, והצוות יעדכן אותם לאחר פגישות עם בעלי העניין
(טבלה~\ref{tab:risks}). לפני הפחתה כל שנים-עשר הסיכונים שזוהו הם אדומים, ולכן לכולם הוגדרה תוכנית הפחתה ומגירה. שלושה מהם נבחנים לעומק משום
שהם נוגעים לליבת הפרויקט: הזיות מודל השפה המקומי (\risk{03}), עיכובי סיווג נתונים בתעשייה האווירית (\risk{01}) וצוואר הבקבוק של זמינות המומחים (\risk{06}).
\begin{table}[!htbp]
\centering
\caption{רשימת סיכונים וכימותם: אדום \mth{R\ge 12}, כתום \rng{5}{11}, ירוק \rng{1}{4}}\label{tab:risks}
\begin{pctable}{c,H{0.52},C{0.11},C{0.11},C{0.12}}
  \pchead{מזהה & תיאור הסיכון & \mth{P} & \mth{I} & \mth{R=P\times I}}
  \riskrow{\risk{01}}{נתונים: עיכוב בסיווג ובאישור הקורפוס הארגוני בתעשייה האווירית}{4}{4}
  \riskrow{\risk{02}}{אבטחה: דליפת מידע או תקשורת לרשת חיצונית}{3}{5}
  \riskrow{\risk{03}}{איכות: הזיית מודל השפה המקומי, תוכן שאינו נתמך במקור}{4}{5}
  \riskrow{\risk{04}}{לוח זמנים: אי-עמידה בלוח הסמסטר}{4}{4}
  \riskrow{\risk{05}}{טכני: איכות אחזור נמוכה בטבלאות}{4}{3}
  \riskrow{\risk{06}}{ארגוני: צוואר בקבוק בזמינות מומחה לאישור}{3}{4}
  \riskrow{\risk{07}}{ארגוני: היענות נמוכה של עובדים ותיקים לראיונות}{3}{4}
  \riskrow{\risk{08}}{תשתית: תלות בתשתית חיצונית}{3}{4}
  \riskrow{\risk{09}}{מחקרי: מדגם ניסוי קטן מדי}{4}{3}
  \riskrow{\risk{10}}{נתונים: תחזוקת גרסאות לא תקינה}{4}{4}
  \riskrow{\risk{11}}{בטיחות: סצנת \latin{3D} מציגה תהליך שגוי או לא מאושר}{4}{4}
  \riskrow{\risk{12}}{משתמש: תקלת נגישות, ביצועים או עומס למשתתף}{3}{4}
\end{pctable}
\end{table}

\needspace{4\baselineskip}
\section{מטריצת סיכונים}
איור~\ref{fig:riskmatrix} ממקם את הסיכונים במטריצה. הריכוז הגבוה ביותר הוא בתאים \mth{(P,I)=(4,4)} ו-\mth{(3,4)}, ו-\risk{03}
(תוכן שאינו נתמך במקור, \mth{R=20}) הוא החמור ביותר --- בהתאם לליקויים שזיהה מאמר הדגל \cite{hua2025civask}.
\begin{figure}[!htbp]
\centering
\pcFig{risk_matrix}{%
\pcAutoFit{\begin{pctikz}[x=1.6cm,y=0.95cm]
  \foreach \p in {1,...,5}{ \foreach \i in {1,...,5}{
    \pgfmathtruncatemacro{\rr}{\p*\i}
    \ifnum\rr>11 \def\cc{pcBad!55}\else\ifnum\rr>4 \def\cc{pcWarn!55}\else\def\cc{pcGood!45}\fi\fi
    \fill[\cc,draw=white,line width=1pt] ({\i-1},{\p-1}) rectangle (\i,\p);
  }}
  \foreach \k in {1,...,5}{
    \node[font=\scriptsize\sffamily] at ({\k-0.5},-0.4) {\k};
    \node[font=\scriptsize\sffamily] at (-0.4,{\k-0.5}) {\k};
  }
  \node[font=\small\sffamily] at (2.5,-1.0) {\tn{השפעה \latin{(I)}}};
  \node[font=\small\sffamily,rotate=90] at (-1.0,2.5) {\tn{הסתברות \latin{(P)}}};
  \pcRiskDots
  % Hebrew zone legend
  \foreach \zc/\zt/\zy in {pcBad!55/{אדום: \mth{R\ge12} --- הפחתה ומגירה}/-1.6, pcWarn!55/{כתום: \rng{5}{11} --- מעקב}/-2.1, pcGood!45/{ירוק: \mth{R\le4} --- קבלה}/-2.6}{
    \fill[\zc] (1.5,{\zy-0.15}) rectangle (1.7,{\zy+0.15});
    \node[anchor=west,font=\scriptsize\sffamily] at (1.75,\zy) {\tn{\zt}};
  }
\end{pctikz}}}
\caption{מטריצת סיכונים \latin{5}$\times$\latin{5}: מיקום כל סיכון לפי הסתברות והשפעה, על בסיס הנחות תכן}\label{fig:riskmatrix}
\end{figure}

\needspace{4\baselineskip}
\section{תוכנית הפחתה, תוכנית מגירה וסיכון שיורי}\label{sec:mitigation}
לכל סיכון אדום מוגדרים מניעה, טריגר מדיד ופעולה בעת כשל (טבלה~\ref{tab:mitigation}). הסיכון השיורי מחושב מחדש אחרי ההפחתה לפי
\begin{equation}
  R_{residual}=P_{new}\times I_{new},
\end{equation}
כאשר \mth{P_{new}} ו-\mth{I_{new}} הם ההסתברות וההשפעה המשוערות לאחר יישום ההפחתה (בטבלה: \mth{P'} ו-\mth{I'}). האחראי לטיפול בכל סיכון
הוא המבצע \latin{(R)} במשימה הקשורה בטבלה~\ref{tab:raci}, והמנחה האקדמי \latin{(A)} מאשר את התוכנית. ציוני השיורי הם הערכת תכן שתיבדק מחדש
לאחר ביצוע ההפחתה, ואינם תוצאות מדידה.
\begin{table}[!htbp]
\centering
\caption{תוכנית הפחתה ותוכנית מגירה לסיכונים עם \mth{R\ge 12}, כולל סיכון שיורי \mth{R_{residual}=P'\times I'}}\label{tab:mitigation}
\renewcommand{\pcTabFont}{\footnotesize}
\newcommand{\resrow}[6]{\pcrow{#1 & #2 & #3 & #4 & \num{#5} & \num{#6} & \textbf{\num{\number\numexpr#5*#6\relax}}}}
\begin{pctable}{c,H{0.19},H{0.30},H{0.14},C{0.06},C{0.06},C{0.08}}
  \pchead{מזהה & הפחתה (מניעה) & תוכנית מגירה (טריגר ופעולה) & אחראי ומועד & \mth{P'} & \mth{I'} & \mth{R_{res}}}
  \resrow{\risk{01}}{בקשת אישור מוקדמת; קורפוס פומבי חלופי מוכן מראש}{טריגר: אין אישור סיווג עד סוף שבוע \latin{2}. פעולה: מעבר לקורפוס \latin{FAA} בלבד}{\StudentTwo, שבוע \latin{2}}{2}{3}
  \resrow{\risk{02}}{בידוד רשת, סינון לפי הרשאה ובדיקת דליפה \cite{zeng2024good}}{טריגר: ניסיון תקשורת חוץ או דליפה בבדיקה. פעולה: עצירת הריצה ובדיקה חוזרת}{\StudentOne, כל ריצה}{1}{5}
  \resrow{\risk{03}}{עיגון בגרסה, הימנעות ממענה, אישור מומחה ושיפוט כפול במדגם}{טריגר: נאמנות נמוכה מ-\pct{90}. פעולה: הקפאת פרסום ומעבר למצב שליפה בלבד, ללא יצירה}{\StudentTwo, \latin{M4}}{2}{4}
  \resrow{\risk{04}}{אבני דרך ושבועות רזרבה \latin{B1}, \latin{B2}}{טריגר: אין אב-טיפוס בסוף שבוע \latin{9}. פעולה: צמצום לשלושה צמתים גנריים}{\StudentOne, שבוע \latin{9}}{3}{2}
  \resrow{\risk{05}}{סימון טבלאות ובדיקת אחזור נפרדת}{טריגר: הצלחה נמוכה מ-\pct{80} בשאלות טבלה. פעולה: הפניה ידנית למקור}{\StudentOne, \latin{M4}}{2}{3}
  \resrow{\risk{06}}{חלונות ביקורת קבועים; תקרת קצב \latin{10} יחידות בשבוע לפי מודל \latin{M/M/1}}{טריגר: המתנה ממוצעת מעל חצי שבוע או פחות משתי ביקורות עד שבוע \latin{10}. פעולה: הקטנת סט השאלות ובקשת מבקר נוסף}{\StudentTwo, שבוע \latin{10}}{2}{3}
  \resrow{\risk{07}}{הסכמה, ראיונות קצרים והיעדר לחץ}{טריגר: פחות משלושה מרואיינים עד שבוע \latin{4}. פעולה: שימוש בתרחישים פומביים}{\StudentTwo, שבוע \latin{4}}{2}{3}
  \resrow{\risk{08}}{מודלים וקורפוס פועלים מקומית}{טריגר: תלות ברשת חיצונית בבדיקת שבוע \latin{5}. פעולה: מעבר למודל מקומי קטן יותר}{\StudentOne, שבוע \latin{5}}{2}{2}
  \resrow{\risk{09}}{מדגם מוכן וגיבוי שאלות}{טריגר: פחות מ-\latin{20} שאלות אמת בשבוע \latin{9}. פעולה: סימון התוצאות כגישושיות}{\StudentTwo, שבוע \latin{9}}{3}{2}
  \resrow{\risk{10}}{\texttt{\latin{document\_id}}, גרסה והשהיית יחידות תלויות (\req{11})}{טריגר: פריט ישן הופיע לאחר עדכון. פעולה: הקפאת פרסום ובנייה מחדש של האינדקס}{\StudentOne, כל עדכון}{2}{3}
  \resrow{\risk{11}}{מומחה מאשר צמתים ומקורות לפני פרסום}{טריגר: פער בין תסריט לסעיף מקור. פעולה: חסימת הסצנה ותיקון מלא לפני שימוש}{\StudentTwo, לפני שימוש}{2}{3}
  \resrow{\risk{12}}{בדיקה מוקדמת בדפדפן והנגשה \mbox{דו-ממדית}}{טריגר: פחות מ-\pct{80} משלימים ללא עזרה. פעולה: מסלול \latin{2D} ותיקון הממשק}{\StudentOne, \latin{M4}}{2}{2}
\end{pctable}
\end{table}
\textbf{שלושת הסיכונים הקריטיים לליבת הפרויקט.} \textbf{הזיות מודל שפה מקומי (\risk{03}):} מודל בגודל \latin{7B}--\latin{8B} בדיוק נמוך מועד יותר להזיה
ממודל גדול; לכן אין סומכים על המודל לבדו אלא על שלוש שכבות בקרה (עיגון והימנעות, אישור מומחה, שיפוט כפול), שהפחיתו במודל התכן את השיעור
הצפוי של טענות לא נתמכות ל-\pct{2.4} (פרק~\ref{ch:alt}); אם הנאמנות במדגם תרד מ-\pct{90}, המערכת עוברת למצב שליפה בלבד שבו מוצגים קטעי המקור ללא ניסוח חדש.
\textbf{עיכובי סיווג נתונים ב-\latin{IAI} (\risk{01}):} בקשת הגישה תוגש כבר בפעולה המיידית הראשונה (פרק~\ref{ch:summary}); ללא אישור עד סוף שבוע \latin{2} הפיילוט
ממשיך על קורפוס \latin{FAA} פומבי, ולכן העיכוב אינו עוצר את הנתיב הקריטי. \textbf{צוואר בקבוק של מומחים (\risk{06}):} ההמתנה לאישור היא פונקציה חדה
של קצב הבדיקה (פרק~\ref{ch:alt}: ב-\mth{\mu=14} יחידות בשבוע \mth{W_q=0.18} שבוע, וב-\mth{\mu=11} כבר \mth{0.91}), ולכן נקבעו תקרת קצב הפקה, חלונות ביקורת קבועים מראש ושיעור
ניצול \mth{\rho} שאינו עולה על \latin{0.83}. לאחר ההפחתה אין סיכון בטווח האדום: הגבוה ביותר הוא \risk{03} (\mth{R_{res}=8}), וכולם כתומים או ירוקים.

\textbf{כלל להפעלת תרגיל עומס:} תזמון שלא נמסר מראש מותר בפיילוט אך ורק לאחר הסכמה מוקדמת, בהודעת “תרגיל” שאינה עלולה
להתפרש כאירוע חירום ממשי, ועם אפשרות יציאה ללא ענישה \cite{nist80084}. בכל ספק סביר לנזק או למצב לחץ חריג יבוטל התרגיל
ויינתן תרגול חלופי. כל שינוי בנוהל בטיחות, סיווג או הרשאה מחייב אישור גורם מוסמך, והערכת השיורי אינה מהווה אישור אבטחה או
בטיחות.
\pcChapterEnd{risk}
